\begin{abstract}
High-resolution document VLMs typically encode thousands of visual tokens per page, yet the content of a document page is highly non-uniform: text, formulas, and tables carry most of the semantic information, while blank margins and background regions occupy large portions of the page. Existing compression strategies---fixed downsampling, token pruning, and fixed-ratio merging---either overlook this density heterogeneity or discard tokens outright, making it hard to reach high compression rates without degrading fine-grained OCR performance. We propose BlackHoleMerge, a training-based, content-adaptive method for visual-token compression in document VLMs. Our approach stems from a simple observation: a pretrained document VLM already embeds an annotation-free density signal within its own internal representations---the L2 norm of its ViT Key (K) vectors---which reflects the text and structural density of each page region. BlackHoleMerge translates this intrinsic signal directly into compression decisions inside the visual encoder. It preserves text- and structure-dense regions at fine granularity, while folding blank and redundant regions into these informative tokens so that information is compressed rather than discarded. We then train the model to adapt to the resulting shortened visual sequences. As a result, the model dynamically reallocates its visual budget according to document density; requiring neither external detectors nor predefined compression ratios, our scheme performs compression in a fully adaptive manner. We validate BlackHoleMerge on multiple end-to-end OCR backbones. On OmniDocBench v1.6, it achieves over 70\% average visual-token compression while maintaining nearly lossless parsing accuracy relative to the uncompressed models. These results demonstrate that content-adaptive token compression is an effective route to improving the efficiency of end-to-end document VLMs.11
\end{abstract}